
\chapter{Introdução}
\label{cap:introducao}

\section{Contextualização}
\label{sec:contextualizacao}

O monitoramento da qualidade da água é uma atividade importante para a gestão dos recursos hídricos, a identificação de alterações ambientais e o apoio à tomada de decisão. A disponibilidade de medições obtidas por sensores permite acompanhar parâmetros físico-químicos ao longo do tempo e observar mudanças nas condições dos corpos d'água. Com a ampliação de sistemas automatizados de monitoramento, torna-se possível obter séries temporais com maior frequência de coleta, aumentando o volume de informações que precisam ser armazenadas, processadas e interpretadas.

Apesar dos benefícios desse tipo de monitoramento, a utilização de dados de sensores apresenta desafios. As medições podem conter valores ausentes, ruídos, inconsistências e alterações decorrentes de falhas instrumentais ou de condições ambientais. Além disso, uma variação significativa em determinado parâmetro não representa necessariamente um problema: ela pode decorrer tanto de uma falha de medição quanto de um evento ambiental legítimo. A distinção entre essas situações exige considerar o comportamento temporal dos dados e o contexto em que foram coletados \cite{leigh2019framework}.

Métodos computacionais de análise de dados e inteligência artificial podem contribuir para o processamento dessas informações, por meio da identificação de padrões, da detecção de anomalias e da geração de alertas. Entretanto, a identificação automática de um comportamento incomum não resolve, por si só, todas as necessidades de interpretação. Em aplicações de monitoramento, também pode ser necessário relacionar as medições observadas a informações contextuais, explicar os resultados em linguagem acessível e fornecer elementos que auxiliem especialistas e gestores a decidir quais ações devem ser consideradas.

Nesse cenário, os modelos de linguagem de grande porte (\textit{Large Language Models} -- LLMs) apresentam possibilidades para a interação em linguagem natural, a síntese de informações e a organização de respostas a partir de diferentes fontes de conhecimento. Quando integrados a ferramentas computacionais, bases de dados e mecanismos de recuperação de informação, esses modelos podem fazer parte de sistemas capazes de executar sequências de tarefas, consultar dados e apoiar processos de análise. A geração aumentada por recuperação (\textit{Retrieval-Augmented Generation} -- RAG), por exemplo, permite combinar a geração textual com informações recuperadas de fontes externas \cite{lewis2020rag}.

A incorporação de LLMs em agentes inteligentes amplia essa possibilidade ao permitir que o modelo participe de um ciclo de interpretação do contexto, planejamento, utilização de ferramentas e análise dos resultados obtidos. Dependendo da arquitetura adotada, o sistema pode consultar dados, acionar procedimentos analíticos e organizar as evidências disponíveis para responder a uma solicitação. Contudo, a utilização de um LLM não garante, isoladamente, a correção das informações geradas. A confiabilidade do sistema depende também da qualidade dos dados, dos mecanismos de integração, da verificação das respostas e da possibilidade de supervisão humana.

Diante desse contexto, esta pesquisa investiga o uso de agentes inteligentes baseados em LLMs como mecanismo de apoio à tomada de decisão em sistemas de monitoramento da qualidade da água. O interesse está na integração entre dados de monitoramento, recursos de análise computacional e capacidades de interação em linguagem natural, considerando tanto as possibilidades de aplicação quanto as limitações técnicas e metodológicas envolvidas.

\section{Motivação e Justificativa}
\label{sec:justificativa}

A motivação desta pesquisa decorre da necessidade de transformar grandes volumes de medições de monitoramento em informações úteis para a compreensão de situações e para o apoio à tomada de decisão. A existência de sensores e de mecanismos automatizados de coleta não elimina a necessidade de interpretar os dados, verificar a consistência das observações e compreender os fatores que podem explicar alterações nos parâmetros monitorados.

A análise de séries temporais de qualidade da água envolve desafios que vão além da identificação de valores atípicos. Uma alteração observada pode estar associada a um evento ambiental real, a uma falha do sensor ou a um problema de transmissão. Assim, a interpretação deve considerar o comportamento dos dados, a informação contextual disponível e os limites dos métodos utilizados. Abordagens automatizadas podem apoiar essa tarefa, mas seus resultados precisam ser apresentados de forma compreensível e passível de verificação.

Os LLMs oferecem uma alternativa para facilitar a interação com sistemas computacionais por meio da linguagem natural e para sintetizar informações provenientes de diferentes fontes. Quando associados a agentes e ferramentas externas, podem atuar em fluxos de trabalho que envolvem consultas, análises e organização de resultados. Essa combinação justifica investigar arquiteturas que não utilizem o modelo de linguagem como única fonte de decisão, mas que o integrem a dados e procedimentos computacionais verificáveis.

A relevância científica do trabalho está em investigar como essas capacidades podem ser organizadas em uma arquitetura voltada ao monitoramento da qualidade da água, identificando possibilidades, requisitos e limitações. A pesquisa também busca distinguir o potencial conceitual dessas tecnologias das evidências efetivamente disponíveis sobre sua aplicação e avaliação.

Do ponto de vista prático, espera-se que a investigação contribua para orientar o desenvolvimento de sistemas capazes de auxiliar usuários na consulta e interpretação de informações de monitoramento. A proposta não pressupõe substituir especialistas nem transferir decisões críticas integralmente para um modelo de linguagem. O foco é estudar mecanismos computacionais que possam apoiar a análise, explicitar informações relevantes e manter a participação humana no processo decisório.

\section{Problema de pesquisa}
\label{sec:problema}

Sistemas de monitoramento da qualidade da água podem produzir dados em alta frequência, cuja interpretação exige considerar o comportamento temporal das medições, a possibilidade de falhas instrumentais e o contexto ambiental. Embora métodos de análise de dados possam identificar padrões e anomalias, os resultados precisam ser relacionados a informações contextuais e apresentados de modo que auxiliem a compreensão das situações observadas.

Os agentes inteligentes baseados em LLMs apresentam recursos que podem contribuir para essa integração, especialmente quando associados a ferramentas analíticas, fontes externas de conhecimento e mecanismos de consulta a dados. Entretanto, ainda é necessário compreender como organizar esses componentes de maneira adequada ao domínio, quais tarefas podem ser apoiadas e quais limitações devem ser consideradas para evitar respostas sem fundamentação suficiente.

Dessa forma, o problema de pesquisa consiste em investigar como agentes inteligentes baseados em LLMs podem ser integrados a dados e métodos de análise de sistemas de monitoramento da qualidade da água para apoiar a interpretação das informações e a tomada de decisão, considerando a confiabilidade, a rastreabilidade e a participação humana.

\section{Questões de pesquisa}
\label{sec:questoes-pesquisa}

A pesquisa é orientada pelas seguintes questões:

\begin{itemize}
    \item \textbf{QP1 -- Aplicações de LLMs e IA generativa:} Quais aplicações de LLMs, IA generativa, geração aumentada por recuperação (RAG) ou agentes inteligentes têm sido propostas para o monitoramento, a interpretação ou a gestão da qualidade da água e de sistemas hídricos?

    \item \textbf{QP2 -- Dados de sensores e monitoramento:} Quais tipos de dados de sensores, Internet das Coisas (IoT) ou redes de monitoramento contínuo são utilizados, e como são tratados, analisados ou interpretados?

    \item \textbf{QP3 -- Mecanismos de apoio à decisão:} Quais mecanismos de apoio à decisão são descritos, incluindo alertas, diagnóstico de anomalias, explicações, recomendações, visualizações e interfaces entre pessoas e sistemas de inteligência artificial?

    \item \textbf{QP4 -- Limitações e lacunas:} Quais limitações técnicas, metodológicas e de avaliação são relatadas, e quais oportunidades de pesquisa permanecem em aberto?
\end{itemize}

Essas questões orientam a análise do estado da arte e contribuem para identificar requisitos relevantes para a definição da proposta de pesquisa.

\section{Objetivos}
\label{sec:objetivos}

\subsection{Objetivo geral}
\label{subsec:objetivo-geral}

Investigar e propor uma abordagem baseada em agentes inteligentes com LLMs para apoiar a interpretação de dados e a tomada de decisão em sistemas de monitoramento da qualidade da água, considerando a integração com fontes de dados e ferramentas analíticas, a confiabilidade das informações produzidas e a supervisão humana.

\subsection{Objetivos específicos}
\label{subsec:objetivos-especificos}

Para alcançar o objetivo geral, são definidos os seguintes objetivos específicos:

\begin{itemize}
    \item Sistematizar os fundamentos teóricos relacionados a agentes inteligentes, LLMs, RAG, dados de sensores, monitoramento da qualidade da água e detecção de anomalias.

    \item Analisar criticamente o estado da arte sobre aplicações de LLMs, IA generativa e agentes inteligentes em sistemas hídricos, considerando dados utilizados, mecanismos de apoio à decisão e limitações relatadas.

    \item Identificar requisitos funcionais e aspectos de confiabilidade relevantes para a integração entre agentes baseados em LLMs, dados de monitoramento e ferramentas analíticas.

    \item Definir uma arquitetura ou abordagem conceitual para apoiar a consulta, a interpretação e a contextualização de informações relacionadas ao monitoramento da qualidade da água.

    \item Estabelecer procedimentos e critérios para avaliar a abordagem proposta, considerando sua adequação às tarefas definidas, a fundamentação das respostas e as limitações identificadas.
\end{itemize}

\section{Método de pesquisa}
\label{sec:metodo-pesquisa}

A pesquisa está organizada em etapas complementares, que articulam a fundamentação teórica, a análise do estado da arte, a definição da proposta e seu planejamento de avaliação.

Inicialmente, será realizada a consolidação da fundamentação teórica sobre inteligência artificial, agentes inteligentes, LLMs, mecanismos de recuperação de informação, dados de sensores e detecção de anomalias. Essa etapa fornece os conceitos necessários para delimitar o problema e compreender os componentes tecnológicos que podem integrar a abordagem proposta.

Em seguida, será analisado o estado da arte por meio de uma revisão de escopo da literatura. O objetivo dessa revisão é mapear aplicações, tipos de dados, mecanismos de apoio à decisão e limitações relacionadas ao uso de IA, LLMs e agentes inteligentes no contexto hídrico. O protocolo e o relato da condução da revisão são apresentados nos apêndices, enquanto o Capítulo~\ref{cap:estado-da-arte} será dedicado à análise comparativa e qualitativa do conjunto de estudos.

A partir da fundamentação e da análise da literatura, serão identificados requisitos para a abordagem proposta. Essa etapa deverá considerar as fontes de dados disponíveis, as tarefas analíticas que podem ser executadas por ferramentas externas, o papel do LLM, os mecanismos de recuperação de informação e as formas de apresentar resultados ao usuário. A definição final da arquitetura dependerá da análise dos requisitos e das condições de acesso aos dados.

Posteriormente, será especificada a abordagem e, conforme a viabilidade técnica e a disponibilidade de dados, será desenvolvido um protótipo ou artefato computacional para investigar sua aplicação. A avaliação deverá ser definida de acordo com as funcionalidades efetivamente implementadas e com os dados disponíveis, evitando atribuir ao sistema capacidades que não tenham sido verificadas experimentalmente.

Por fim, os resultados serão analisados em relação aos objetivos da pesquisa, às questões identificadas no estado da arte e às limitações do método. Os procedimentos específicos de implementação e avaliação serão detalhados no capítulo metodológico correspondente.

\section{Cronograma proposto}
\label{sec:cronograma}

O cronograma a seguir apresenta uma previsão inicial das atividades. Os períodos deverão ser ajustados conforme o calendário do programa, as orientações do orientador e a disponibilidade dos dados necessários à avaliação.

\begin{table}[htbp]
    \centering
    \caption{Cronograma proposto para o desenvolvimento da pesquisa}
    \label{tab:cronograma-pesquisa}
    \small
    \begin{tabular}{p{7.5cm}p{2.5cm}}
        \toprule
        \textbf{Atividade} & \textbf{Período previsto} \\
        \midrule
        Fazer direito isso & Out.--Nov./2026 \\
        %Consolidação da fundamentação teórica & Out.--Nov./2026 \\
        %Leitura integral e análise dos estudos selecionados & Out.--Dez./2026 \\
        %Síntese do estado da arte e identificação de lacunas & Nov./2026--Jan./2027 \\
        %Definição dos requisitos e da arquitetura proposta & Jan.--Fev./2027 \\
        %Desenvolvimento ou especificação do protótipo & Fev.--Abr./2027 \\
        %Planejamento e execução da avaliação & Abr.--Mai./2027 \\
        %Análise dos resultados e discussão & Mai.--Jun./2027 \\
        %Redação e revisão dos capítulos finais & Jun.--Jul./2027 \\
        %Revisão geral, ajustes e preparação para defesa & Jul./2027 \\
        \bottomrule
    \end{tabular}
    \fonte{Elaboração própria.}
\end{table}

\section{Limitações e riscos}
\label{sec:limitacoes-riscos}

A pesquisa apresenta limitações relacionadas à disponibilidade e à qualidade dos dados, à maturidade das tecnologias investigadas e à possibilidade de avaliar a abordagem em condições representativas de uso.

Uma primeira limitação diz respeito à disponibilidade de dados de monitoramento. A ausência de séries temporais suficientemente longas, dados rotulados ou informações contextuais pode restringir a avaliação de determinadas funcionalidades, especialmente aquelas relacionadas à detecção de anomalias e à interpretação de eventos ambientais. Caso essas restrições ocorram, o escopo experimental deverá ser ajustado e as consequências para a validade dos resultados serão explicitadas.

Outra limitação está relacionada à revisão de escopo. A cobertura da literatura depende das fontes consultadas, das expressões de busca e dos critérios de seleção adotados. Além disso, a seleção bibliográfica preliminar não equivale à confirmação da elegibilidade dos estudos após leitura integral. A análise do estado da arte deverá considerar essas restrições e distinguir evidências consolidadas de propostas conceituais ou resultados ainda insuficientemente avaliados.

O uso de LLMs também envolve riscos de geração de informações incorretas, respostas sem fundamentação suficiente e interpretações inadequadas de dados. Para reduzir esses riscos, a abordagem deverá priorizar a consulta a fontes identificáveis, a utilização de ferramentas analíticas apropriadas, a explicitação das incertezas e a verificação das informações utilizadas para sustentar as respostas. Sempre que necessário, os resultados deverão ser apresentados como apoio à análise, e não como decisão automática.

Há ainda riscos de integração entre componentes, incluindo incompatibilidades entre formatos de dados, limitações de acesso a serviços externos, custos computacionais e variações no comportamento dos modelos. Esses aspectos deverão ser considerados na escolha da arquitetura e no planejamento da avaliação.

Por fim, a generalização dos resultados poderá ser limitada pelo domínio, pelos conjuntos de dados e pelas tarefas utilizadas nos experimentos. A pesquisa deverá registrar as condições em que a abordagem for avaliada e evitar extrapolar conclusões para contextos que não tenham sido investigados.

\section{Estrutura do documento}
\label{sec:estrutura-documento}

Além deste capítulo introdutório, o documento está organizado em capítulos que desenvolvem os fundamentos, a análise da literatura, o método e a avaliação da pesquisa.

O Capítulo~2 apresenta a fundamentação teórica necessária para compreender os conceitos de inteligência artificial, agentes inteligentes, LLMs, geração aumentada por recuperação, arquiteturas baseadas em agentes, dados de sensores e detecção de anomalias.

O Capítulo~3 apresenta a análise do estado da arte, com uma discussão geral, comparativa e qualitativa dos estudos relacionados ao uso de inteligência artificial, LLMs e agentes inteligentes no monitoramento e na gestão da qualidade da água. O capítulo busca identificar padrões, diferenças entre abordagens, limitações e lacunas relevantes para a pesquisa.

O Capítulo~4 descreve a metodologia da pesquisa e os procedimentos previstos para especificar, desenvolver ou implementar a abordagem proposta e planejar sua avaliação.

O Capítulo~5 apresenta os resultados obtidos com a aplicação ou avaliação da abordagem, incluindo a análise de seu comportamento e a discussão das evidências produzidas.

Por fim, o Capítulo~6 reúne as conclusões, retoma os objetivos da pesquisa, discute as contribuições e limitações do trabalho e apresenta possibilidades de continuidade.
